\chapter{Signal Timing}
\label{ch:signal-timing}

This chapter gives the timing of the pins of the device (Chapter~\ref{ch:package}): the bus cycles in which the CPU
reads and writes it, the interrupt outputs, reset, and the video output. Chapter~\ref{ch:timing} gives the timing
inside the device, in the same units.

\section{Reading the Figures}

The figures follow the conventions of the bus timing diagrams of the \hwmanual{}. Each labeled interval is one
\texttt{CLKI} cycle, which is one PPU cycle, seen between rising edges. Waveform rows show physical levels: for an
active-low signal, \texttt{L} means asserted. Bands show a value; an empty band carries no defined value, and a
three-state signal that is not driven is drawn as a flat line at mid-level. The pulses in the \texttt{Sample} row mark
the rising edges at which the device samples the bus. Electrical delays are not specified, so the figures show logical
timing only.

\section{Bus Cycles}
\label{sec:signal-bus}

The device answers every cycle in which \activelow{CS} and \activelow{BAS} are both asserted, and it answers in the
cycle itself: \activelow{BACK} is asserted from the start of the cycle, so the CPU accepts the cycle at the next rising
edge and never waits (Chapter~\ref{ch:package}).

\subsection{Writes}

The CPU presents the offset on A0--A15, the data on D0--D15 and BRW low, just after the rising edge that starts the
cycle. The device takes the data at the rising edge that ends the cycle; the register, palette entry, video memory word
or sprite table word holds the new value from that edge, and a read in the next cycle returns it. A write to an offset
that is not stored, or one that the rules of Section~\ref{sec:memory-vram-access} discard, is acknowledged in the same
way.

\begin{figure}[H]
	\centering
	\small
	\begin{tikztimingtable}[sirc timing diagram]
		CLKI             & L 2{0.5H 0.5L}              \\
		Phase            & D{access} D{next} D{}       \\
		\activelow{CS}   & L H H                       \\
		\activelow{BAS}  & L H H                       \\
		BRW              & L H H                       \\
		A0--A15          & D{offset} D{} D{}           \\
		D0--D15          & D{data} Z Z                 \\
		\activelow{BACK} & L H H                       \\
		Stored word      & D{old} 2D{new}              \\
		Sample           & 0.95L 0.1H 1.95L            \\
	\end{tikztimingtable}
	\caption{Write Cycle}
	\label{fig:signal-write}
\end{figure}

\subsection{Reads}

On a read the device drives D0--D15 with the word the offset selects, from the start of the cycle to its end, and
releases the bus after it. The word is valid before the rising edge that ends the cycle, which is the edge at which the
CPU samples it. A read changes nothing in the device; no register has a side effect on a read.

\begin{figure}[H]
	\centering
	\small
	\begin{tikztimingtable}[sirc timing diagram]
		CLKI             & L 2{0.5H 0.5L}              \\
		Phase            & D{access} D{next} D{}       \\
		\activelow{CS}   & L H H                       \\
		\activelow{BAS}  & L H H                       \\
		BRW              & H H H                       \\
		A0--A15          & D{offset} D{} D{}           \\
		D0--D15          & D{data} Z Z                 \\
		\activelow{BACK} & L H H                       \\
		Sample           & 0.95L 0.1H 1.95L            \\
	\end{tikztimingtable}
	\caption{Read Cycle}
	\label{fig:signal-read}
\end{figure}

\subsection{Consecutive Cycles}

The device accepts one access in every cycle, with no recovery time between accesses. A DMA transfer, which makes one
bus cycle for each word (\hwmanual{}), therefore moves one word per \texttt{CLKI} cycle into the palette, the sprite
table or video memory, whichever the offsets address.

\begin{figure}[H]
	\centering
	\small
	\begin{tikztimingtable}[sirc timing diagram]
		CLKI             & L 5{0.5H 0.5L}                                        \\
		Phase            & D{Setup} D{Write 1} D{Write 2} D{Write 3} D{} D{}     \\
		\activelow{CS}   & H L L L H H                                           \\
		\activelow{BAS}  & H L L L H H                                           \\
		BRW              & H L L L H H                                           \\
		A0--A15          & D{} D{v} D{v+1} D{v+2} D{} D{}                        \\
		D0--D15          & Z D{d0} D{d1} D{d2} Z Z                               \\
		\activelow{BACK} & H L L L H H                                           \\
		Last word stored & 2D{} D{v} D{v+1} 2D{v+2}                              \\
		Sample           & 1.95L 0.1H 0.9L 0.1H 0.9L 0.1H 1.95L                  \\
	\end{tikztimingtable}
	\caption{Three-Word DMA Write into the Device}
	\label{fig:signal-dma}
\end{figure}

\subsection{The Start of Vertical Blanking}
\label{sec:signal-vblank}

Video memory and the sprite table are open to the CPU from the first cycle of line 224 to the last cycle of line 261,
and at any time while \texttt{GD} is set (Section~\ref{sec:memory-vram-access}). The gate changes at a line boundary,
so it is exact to the cycle. A write in the last cycle of line 223 is discarded; a write in the first cycle of line 224
is stored. The vertical sync interrupt is raised in that same first cycle (Section~\ref{sec:interrupts}), and the CPU
takes it at the end of the instruction in progress, so the handler's first access to video memory is some cycles later
than the figure shows.

\begin{figure}[H]
	\centering
	\small
	\begin{tikztimingtable}[sirc timing diagram]
		CLKI              & L 3{0.5H 0.5L}                     \\
		Line              & D{223} D{223} D{224} D{224}        \\
		Cycle in line     & D{680} D{681} D{0} D{1}            \\
		\activelow{CS}    & L L L L                            \\
		\activelow{BAS}   & L L L L                            \\
		BRW               & L L L L                            \\
		A0--A15           & D{v} D{v+1} D{v+2} D{v+3}          \\
		D0--D15           & D{d0} D{d1} D{d2} D{d3}            \\
		\activelow{BACK}  & L L L L                            \\
		Video memory open & L L H H                            \\
		Word stored       & D{no} D{no} D{yes} D{yes}          \\
		\activelow{NMI}   & H H L H                            \\
	\end{tikztimingtable}
	\caption{Video Memory Writes across the Start of Vertical Blanking (\texttt{VIE} set)}
	\label{fig:signal-vblank}
\end{figure}

\section{Interrupt Timing}
\label{sec:signal-interrupts}

The interrupt outputs change just after a rising edge of \texttt{CLKI}, like the other outputs of the device. The CPU
samples them at the rising edge that ends the cycle.

\subsection{Vertical Sync}

Figure~\ref{fig:signal-vblank} shows the pulse: \activelow{NMI} is asserted for the whole of the first cycle of line
224 and deasserted in the next, when \texttt{VIE} is set in the Interrupt Control register. The pulse is one
\texttt{CLKI} cycle, which the CPU always samples, and the device does not wait for or expect an acknowledgement
(Section~\ref{sec:interrupts}).

\subsection{Line Compare}

\activelow{IRQ4} is asserted at the first cycle of the line whose number is in \texttt{LINE}, when \texttt{LIE} is set,
and is held. It is released by the program, with a write to \texttt{LINE}, or by clearing \texttt{LIE}.
The \texttt{LIP} bit of the Status register follows the output.

\begin{figure}[H]
	\centering
	\small
	\begin{tikztimingtable}[sirc timing diagram]
		CLKI              & L 3{0.5H 0.5L}                     \\
		Line              & D{N-1} D{N-1} D{N} D{N}            \\
		Cycle in line     & D{680} D{681} D{0} D{1}            \\
		\texttt{LINE}     & 4D{N}                              \\
		\texttt{LIE}      & 4H                                 \\
		\activelow{IRQ4}  & H H L L                            \\
		\texttt{LIP}      & L L H H                            \\
	\end{tikztimingtable}
	\caption{Line-Compare Interrupt Raised at the Start of Line N}
	\label{fig:signal-irq-raise}
\end{figure}

The acknowledging write is an ordinary write cycle. The device releases \activelow{IRQ4} and clears \texttt{LIP} at the
rising edge that ends the write, so the line is high in the cycle after it. A program should write \texttt{LINE} early
in its handler, and before it returns, because \activelow{IRQ4} is level-sensitive and the CPU recognises the interrupt
again if the line is still asserted when the handler ends (Section~\ref{sec:interrupts}).

\begin{figure}[H]
	\centering
	\small
	\begin{tikztimingtable}[sirc timing diagram]
		CLKI              & L 3{0.5H 0.5L}                     \\
		Phase             & D{access} D{next} D{} D{}          \\
		\activelow{CS}    & L H H H                            \\
		\activelow{BAS}   & L H H H                            \\
		BRW               & L H H H                            \\
		A0--A15           & D{0x0015} D{} D{} D{}              \\
		D0--D15           & D{line} Z Z Z                      \\
		\activelow{BACK}  & L H H H                            \\
		\texttt{LINE}     & D{old} 3D{line}                    \\
		\activelow{IRQ4}  & L H H H                            \\
		\texttt{LIP}      & H L L L                            \\
		Sample            & 0.95L 0.1H 2.95L                   \\
	\end{tikztimingtable}
	\caption{Acknowledging the Line-Compare Interrupt with a Write to \texttt{LINE}}
	\label{fig:signal-irq-ack}
\end{figure}

\section{Reset}
\label{sec:signal-reset}

\activelow{RST} is sampled at each rising edge of \texttt{CLKI}. Every register takes its reset value at the first edge
at which \activelow{RST} is sampled asserted (Table~\ref{tab:reset}), and the outputs go to their idle state: the
interrupt outputs are deasserted, \activelow{CSYNC} is high, and the colour outputs are at black. The line and cycle
counters are held at the first cycle of line 0 while \activelow{RST} is asserted. The first cycle of line 0 is the cycle
that follows the rising edge at which \activelow{RST} is first sampled deasserted. The contents of the palette, video
memory and sprite table are undefined after reset (Chapter~\ref{ch:registers}).

\begin{figure}[H]
	\centering
	\small
	\begin{tikztimingtable}[sirc timing diagram, timing/xunit=3em]
		CLKI              & L 9{0.5H 0.5L}                                   \\
		\activelow{RST}   & 6L 4H                                            \\
		Registers         & D{any} 9D{reset values}                          \\
		Line and cycle    & D{} 6D{held at line 0, cycle 0} D{0} D{1} D{2}   \\
		Interrupts        & 10H                                           \\
		\activelow{CSYNC} & 7H 3L                                            \\
		R, G, B           & 10D{black}                                       \\
	\end{tikztimingtable}
	\caption{Reset, with \activelow{RST} Held for the Minimum of Six Cycles}
	\label{fig:signal-reset}
\end{figure}

\section{Video Output}
\label{sec:video-output}

The video outputs follow the line timing of Chapter~\ref{ch:timing}. A new pixel appears on R, G and B every four
master clocks, at the pixel boundary, and the colour outputs are at black outside the visible part of the line, in the
vertical blanking interval, and in forced blank. The brightness setting scales the colour before conversion
(Section~\ref{sec:reg-base-config}).

\subsection{Horizontal Sync}

\activelow{CSYNC} is asserted for the first 100 master clocks (50 cycles) of every line, and is high for the remaining 1264.
The pulse begins at the first cycle of the line, the same cycle at which the line-compare interrupt is raised. The back
porch follows it, and the first pixel is 160 master clocks after the end of the pulse. The line ends with the front
porch, 80 master clocks; the 64 clocks before the first pixel and the 44 before the front porch are black, so the picture
has a narrow black border on each side.

\begin{figure}[H]
	\centering
	\begin{tikzpicture}[
			x=0.0057\textwidth,
			y=0.62cm,
			lane/.style={font=\scriptsize, anchor=east},
			marker/.style={dashed, thick, sircgray},
			note/.style={font=\scriptsize, text=sircgray}
		]
		% x is in units of 10 master clocks
		\node[lane] at (-4,3.3) {\activelow{CSYNC}};
		\draw[thick] (-4,3.9) -- (0,3.9) -- (0,3.0) -- (10,3.0) -- (10,3.9) -- (140,3.9);
		\node[lane] at (-4,1.4) {R, G, B};
		\draw[thick] (-4,1.1) -- (26,1.1) -- (26,1.8) -- (128.4,1.8) -- (128.4,1.1) -- (140,1.1);
		\fill[sirclightgray] (26,1.1) rectangle (128.4,1.8);
		\node[font=\scriptsize] at (77.2,1.45) {visible: 256 pixels of 4 clocks};
		\node[note] at (5,2.45) {sync};
		\node[note] at (18,2.45) {back porch, 160};
		\node[note] at (132.4,2.45) {front porch, 80};
		\draw[marker] (0,0.6) -- (0,4.6);
		\draw[marker] (10,0.6) -- (10,4.6);
		\draw[marker] (26,0.6) -- (26,4.6);
		\draw[marker] (128.4,0.6) -- (128.4,4.6);
		\draw[marker] (136.4,0.6) -- (136.4,4.6);
		\foreach \x/\l in {0/0, 10/100, 26/260, 128.4/1284, 136.4/1364} { \node[font=\tiny, anchor=north] at (\x,0.6) {\l}; }
		\node[note, anchor=north] at (68,-0.2) {master clocks from the start of the line};
	\end{tikzpicture}
	\caption{Horizontal Sync and the Visible Part of a Line}
	\label{fig:signal-hsync}
\end{figure}

\subsection{Vertical Sync}

Lines 0 to 223 are the picture. The vertical sync is three lines long and ends 14 lines before the first line of the
picture: lines 245, 246 and 247 on NTSC, and lines 295, 296 and 297 on PAL. In each of those lines \activelow{CSYNC} is
asserted from the first cycle of the line for 1272 master clocks (636 cycles) and is high for the last 92. The other
lines of the vertical blanking interval carry the ordinary horizontal sync pulse, and there are no equalising pulses.

\begin{figure}[H]
	\centering
	\begin{tikzpicture}[
			x=0.0012\textwidth,
			y=0.62cm,
			lane/.style={font=\scriptsize, anchor=east},
			note/.style={font=\scriptsize, text=sircgray}
		]
		% x is in units of 10 master clocks; five lines of 136.4 units
		\node[lane] at (-8,0.5) {\activelow{CSYNC}};
		\foreach \i in {0,1,2,3,4} {
				\pgfmathsetmacro{\o}{\i*136.4}
				\ifnum\i=0 \draw[thick] (\o,1) -- (\o,0) -- (\o+10,0) -- (\o+10,1);
				\else\ifnum\i=4 \draw[thick] (\o,1) -- (\o,0) -- (\o+10,0) -- (\o+10,1);
					\else
						\draw[thick] (\o,1) -- (\o,0) -- (\o+127.2,0) -- (\o+127.2,1) -- (\o+136.4,1);
					\fi\fi
			}
		\draw[thick] (-8,1) -- (0,1);
		\draw[thick] (10,1) -- (136.4,1);
		\draw[thick] (4*136.4+10,1) -- (682+8,1);
		\foreach \i/\l in {0/244, 1/245, 2/246, 3/247, 4/248} {
				\pgfmathsetmacro{\o}{\i*136.4}
				\draw[dashed, sircgray] (\o,-0.3) -- (\o,1.4);
				\node[note, anchor=north] at (\o+68.2,-0.35) {line \l};
			}
	\end{tikzpicture}
	\caption{Vertical Sync in Lines 245 to 247 (NTSC), with Horizontal Sync before and after}
	\label{fig:signal-vsync}
\end{figure}
